\begin{abstract}
Unsupervised anomaly detection in multi-sensor Internet-of-Things (IoT) networks and remote sensing telemetry supports operational health monitoring across aerospace systems, planetary exploration, cyber-physical infrastructure, and wearable sensing. However, detecting operational faults from sensor streams presents key challenges: reconstructive deep models (e.g., Autoencoders, Masked Transformers) often fit high-frequency sensor noise; heuristic contrastive detectors can distort state boundaries; and deterministic predictors collapse to the conditional mean across multi-branching dynamics. To address these limitations under edge computational budgets, we propose DyFlow-JEPA (Dynamical Flow-Matching Joint-Embedding Predictive Architecture), a continuous self-supervised framework for unsupervised anomaly detection in multi-sensor and remote sensing telemetry. DyFlow-JEPA frames anomaly detection as tangent velocity vector field modeling on a regularized latent dynamical manifold. DyFlow-JEPA pairs a causal context encoder with an Exponential Moving Average (EMA) target encoder, training a continuous velocity network via Optimal Transport Conditional Flow Matching (OT-CFM) regularized by non-contrastive variance and spectral covariance penalties to prevent representation collapse without synthetic negative sampling. For edge inference, we evaluate a single-pass Deterministic Midpoint Velocity Discrepancy at $t=0.5$, whitened by an empirical Mahalanobis precision tensor and thresholded via Extreme Value Theory (EVT / SPOT) under the Pickands-Balkema-de Haan theorem. Benchmarking across 1,300 model-channel runs on 11 real-world telemetry datasets comprising 130 continuous channels demonstrates that DyFlow-JEPA outperforms deep baselines (TimesNet, TranAD, Anomaly Transformer, DCdetector, NCAD), achieving a channel-weighted Point-F1 of $0.2389$ (+52.5\% over TimesNet and +56.2\% over TranAD) and an unweighted dataset macro Point-F1 of $0.1921$. Phase-space dynamical analyses confirm that DyFlow-JEPA models nominal multi-sensor physical dynamics as stable limit cycles, providing interpretable state trajectories and suppressing false alarms.
\end{abstract}


\begin{IEEEkeywords}
Multi-Sensor IoT, Remote Sensing Telemetry, Spacecraft Telemetry, Anomaly Detection, Joint-Embedding Predictive Architecture (JEPA), Flow Matching, Optimal Transport, Continuous Dynamical Systems, Wearable Edge Sensing, Extreme Value Theory.
\end{IEEEkeywords}
